\documentclass[11pt,a4paper]{article}
\usepackage[utf8]{inputenc}
\usepackage{amsmath,amssymb,amsthm,mathrsfs}
\usepackage[margin=2.4cm]{geometry}
\usepackage{booktabs}
\usepackage[colorlinks,linkcolor=blue,citecolor=blue,urlcolor=blue]{hyperref}
\usepackage{enumitem}
\usepackage{xcolor}

\newtheorem{claim}{Claim}
\newtheorem{observation}[claim]{Observation}
\theoremstyle{definition}
\newtheorem{finding}{Finding}

\newcommand{\eps}{\varepsilon}
\newcommand{\hs}{\mathfrak{h}}
\newcommand{\Gi}{\Gamma_{\!\infty,\pm}}
\newcommand{\DW}{\mathcal{D}_W}
\newcommand{\Dstd}{\mathcal{D}_{\mathrm{std}}}
\newcommand{\sh}{\hat{s}}
\newcommand{\OS}{\textup{[OS]}}

\title{A critical re-reading of\\
\emph{Conformal field theory from lattice fermions}\\
\normalsize T.~J.~Osborne and A.~Stottmeister, Commun.\ Math.\ Phys.\ \textbf{398}, 219--289 (2023); arXiv:2107.13834}
\author{}
\date{29 July 2026}

\begin{document}
\maketitle
\vspace{-2.2em}

\begin{abstract}
\noindent
This is an independent second pass over the technical core of \OS{} (Sections 3.1--3.2, 4.2.3--4.2.4,
5, 6), carried out without deference to the earlier verification note in
\texttt{osborne-stottmeister2023/}, which recorded ``\textsc{correct}, no logical gaps.''
That verdict is too generous. The main theorems still look true, and I found no
counterexample to any of them. But the paper contains \textbf{one invalid proof step
that recurs in three places}, \textbf{one error estimate that is vacuous (indeed $+\infty$)
as stated}, \textbf{a hierarchy of unstated regularity thresholds} whose numerical values are
much larger than the folklore ``$K\ge 2$ suffices'', \textbf{one remark that does not follow
from the lemma it invokes}, \textbf{a normalisation inconsistency that would change the
central charge to $c=(L/\pi)^2$ if read literally}, and \textbf{two unproved ingredients that
carry the whole proof of Theorem 6.3}. All numerical claims below were verified against
explicitly constructed Daubechies filters; the scripts are in \texttt{numerics/}.
\end{abstract}

\section*{Summary of findings}
\begin{center}
\footnotesize
\begin{tabular}{cp{4.1cm}p{6.5cm}l}
\toprule
\# & Location & Issue & Severity \\
\midrule
1 & Lemma 4.5, p.~46 & dominating function: the $\sh$-ratio bound is false & \textbf{gap, repairable} \\
2 & (204)--(210), p.~53--54 & $\mathrm{Error}^2(\delta,L,k,N)=+\infty$ for $k\neq0$ (Ramond) & \textbf{not established} \\
3 & ``sufficiently regular'' & three thresholds, none stated; $\delta{=}2$ needs $K\ge9$ & \textbf{hypothesis gap} \\
4 & Remark 4.14, p.~50 & $P^\pm$ does not preserve $\DW$; claim does not follow & \textbf{gap, repairable} \\
5 & Def.~4.2 (181)/(199) vs (165)/(167) & $\ell_{\pm,k}$ off by $L/\pi$; (170) would give $c=(L/\pi)^2$ & normalisation \\
6 & Thm.~6.3 proof, p.~64 & the two decisive ingredients are asserted, not proved & \textbf{rigor debt} \\
7 & Thm.~4.16 preamble, p.~50 & ``preserving localization (Prop.~3.6)'' not established & overstatement \\
8 & (25), Lem.~3.7, (204)/(205) & typos / vacuous constants & cosmetic \\
\bottomrule
\end{tabular}
\end{center}

\medskip\noindent
Findings 1, 2 and 3 are one phenomenon seen from three sides, so I treat them together in
Section~\ref{sec:ratio}--\ref{sec:reg}. Nothing here contradicts Theorems A, B, C.

\section{The $\hat s$-ratio: a systematic invalid step}\label{sec:ratio}

\subsection{The zeros of $\hat s$}

Throughout, $s$ is a compactly supported orthonormal Daubechies scaling function of order
$K$ (so $2K$ taps), $m_0$ its low-pass symbol and
$\sh(\xi)=\prod_{j\ge1}m_0(2^{-j}\xi)$ as in (102). Two elementary facts are used
nowhere in \OS{} but are decisive:

\begin{observation}\label{obs:zeros}
Orthonormality gives $\sum_{n\in\mathbb Z}|\sh(\xi+2\pi n)|^2=1$; at $\xi=0$ together with
$\sh(0)=1$ this forces $\sh(2\pi n)=0$ for every $n\neq0$. Moreover
$\sh(\xi)=\big(e^{-i\xi/2}\tfrac{\sin(\xi/2)}{\xi/2}\big)^{K}\prod_{j\ge1}L(2^{-j}\xi)$, so the
zero at $2\pi n$ has order exactly $K$.
\end{observation}

\noindent
Numerically (\texttt{numerics/os\_check.py}, block [0]--[1]): $|\sh(2\pi n)| \le 10^{-12}$ for
$\mathrm{db}2,\dots,\mathrm{db}6$ and $n=1,2,3,5$, and
$|\sh(2\pi+\delta)|\sim\delta^{p}$ with $p=2.00,\,2.18,\,3.08$ for $K=2,3,4$.

\subsection{Finding 1: the dominating function in Lemma 4.5}

\begin{finding}\label{f:lemma45}
The domination step in the proof of Lemma 4.5 is invalid, and the stated inequality is false.
\end{finding}

\noindent
The proof reduces convergence to $\ell^2(\Gi,(2L)^{-1})$-convergence of
\[
  f^{(N)}_k(M,m)=\sh(\eps_N m)\,\cos(\tfrac14\eps_Nk)^2\,
  \frac{\sin(\eps_N(m+\tfrac k2))}{\eps_N}\prod_{j=1}^{N-M}m_0\big(\eps_{M+j}(m+k)\big)
  \;\longrightarrow\; f_k(M,m)=\sh(\eps_M(m+k))\,(m+\tfrac k2),
\]
and asserts the existence of an $\ell^2$-dominating function via
\begin{equation}\label{eq:badbound}
  |f^{(N)}_k(M,m)|\;\le\;\Big|(m+\tfrac k2)\sh(\eps_M(m+k))\Big|
  \left|\frac{\sh(\eps_N m)}{\sh(\eps_N(m+k))}\right|
  \;\le\; C_k\Big|(m+\tfrac k2)\sh(\eps_M(m+k))\Big|,
\end{equation}
``for some $C_k>0$, because of the uniform continuity of $\sh$.''

Two things go wrong. First, \emph{uniform continuity of $\sh$ controls the difference
$\sh(\eps_Nm)-\sh(\eps_N(m+k))$, never the quotient}; bounding a quotient needs a lower bound
on $|\sh|$, which is exactly what Observation~\ref{obs:zeros} denies. Second, and concretely:
in the Ramond sector $\Gamma_{\infty,+}=\frac{\pi}{L}\mathbb Z$ one may choose
\[
  m+k=\frac{2\pi n}{\eps_N}=\frac{\pi}{L}\,2^{N+1}n\ \in\ \Gamma_{\infty,+},\qquad n\neq0,
\]
so that $\sh(\eps_N(m+k))=\sh(2\pi n)=0$ while $\sh(\eps_N m)=\sh(2\pi n-\eps_Nk)\neq0$ for
$k\neq0$. The middle expression in \eqref{eq:badbound} is then $\infty$, and the right-hand
side is $0$ (because $\sh(\eps_M(m+k))=\sh(2^{N-M}\!\cdot\!2\pi n)=0$ as well), while the
left-hand side is strictly positive. \emph{No constant $C_k$ makes \eqref{eq:badbound} true.}

\medskip\noindent
Numerically (\texttt{os\_check2.py}, block [4]), for $\mathrm{db}4$, $M=2$, $L=\pi$:

\begin{center}\small
\begin{tabular}{ccrll}
\toprule
$N$ & $k$ & $m$ & $|f^{(N)}_k(M,m)|$ & claimed bound $C_k|f_k(M,m)|$\\
\midrule
6 & 2 & $126$    & $1.547\cdot10^{-6}$  & $C_k\cdot 7.7\cdot10^{-27}$ (i.e.\ $0$)\\
8 & 2 & $510$    & $5.936\cdot10^{-9}$  & $C_k\cdot 3.1\cdot10^{-26}$ (i.e.\ $0$)\\
10& 2 & $2046$   & $2.307\cdot10^{-11}$ & $C_k\cdot 1.2\cdot10^{-25}$ (i.e.\ $0$)\\
\bottomrule
\end{tabular}
\end{center}
(The right-hand column is machine zero: $\sh$ vanishes there to order $K$.)

\subsection{Why the lemma is nevertheless true, and how to repair the proof}

The saving grace is that at these resonances the \emph{sine} factor is also near a zero:
$\eps_N(m+\tfrac k2)=2\pi n-\tfrac12\eps_Nk$, so
$\sin(\eps_N(m+\tfrac k2))/\eps_N\approx -k/2$ rather than $\approx m$. Together with
$|\sh(2\pi n-\eps_Nk)|\asymp c_n(\eps_N|k|)^K$ this gives
$|f^{(N)}_k(M,m)|=O(\eps_N^{\,K})$ at the resonance points. Measured decay
(\texttt{os\_check2.py}, block [6]): $2^{-pN}$ with $p=2.003,\,3.003,\,4.005$ for
$\mathrm{db}2,3,4$ --- i.e.\ exactly $p=K$.

\begin{claim}[repair]
Do not form the quotient. Bound instead
\[
  |f^{(N)}_k(M,m)|\;\le\;\Big|m+\tfrac k2\Big|\;\big|\sh(\eps_N m)\big|
  \prod_{j=1}^{N-M}\big|m_0(\eps_{M+j}(m+k))\big| ,
\]
and estimate the \emph{finite product} directly by the paper's own Lemma 3.7 inside the
$N$-th Brillouin zone ($|m+k|\le\pi/\eps_N$), and by $|\sh(\eps_Nm)|\le C(1+\eps_N|m|)^{-(K-\bar K)}$
together with the $2^{N-M+1}\pi$-periodicity of the product outside it. This yields
$g_k(M,m)=C_{k,M}(1+|m|)^{-(K-\bar K-1)}$, which lies in $\ell^2$ under the same condition as
the paper's own $g_k$.
\end{claim}

\noindent
Lemma 3.7 is proved in \OS{} for precisely this purpose (``\emph{a decay estimate on the finite
products \dots that directly implies (102) \dots by dominated convergence}''), and is then not
used in the proof of Lemma 4.5. The repair is therefore book-keeping, not new mathematics.

\medskip\noindent
Two numerical confirmations (\texttt{os\_check2.py} [5], \texttt{os\_check3.py} [9]):
\begin{itemize}[leftmargin=1.4em,itemsep=1pt]
\item $\sum_m|f^{(N)}_k(M,m)-f_k(M,m)|^2\to0$ monotonically; for $\mathrm{db}4$, $M=2$, $k=2$ it
falls from $1.36\cdot10^{1}$ at $N=4$ to $1.01\cdot10^{-4}$ at $N=14$, with a stable ratio
$\approx0.33$ per step. \emph{The conclusion of Lemma 4.5 is correct.}
\item $G(m):=\sup_{N>M}|f^{(N)}_k(M,m)|$ has the \emph{same} decay exponent as $|f_k(M,m)|$ on
every dyadic window up to $m=2^{13}$, and $\sum_mG(m)^2/\sum_m|f_k(M,m)|^2=1.008$. A valid
dominating function exists; it is just not the one exhibited.
\end{itemize}

\subsection{Finding 2: the error estimates (204)--(210) inherit the singularity}

\begin{finding}\label{f:error}
$\mathrm{Error}^2(\delta,L,k,N)$ as \emph{defined} in (206) equals $+\infty$ for every $k\neq0$
and every $N$ in the Ramond sector. The passage from (208), line 1, to (208), line 2, is
vacuous for $k\neq0$ in both sectors. Consequently (210) --- and with it (207) --- is
established only for $k=0$.
\end{finding}

\noindent
The factorisation (204)--(205) pulls $|(m+\tfrac k2)\sh(\eps_M(m+k))|^2$ out of the sum and
leaves behind the same quotient $\sh(\eps_Nm)/\sh(\eps_N(m+k))$, now inside a supremum. In the
Ramond sector the supremum runs over a set containing an exact pole (same $m$ as above), so
$\mathrm{Error}^2=+\infty$. In (208), line 2, the discrete supremum over $m\in\Gi$ is relaxed to
$\sup_{x\in\mathbb R}$; since $x\mapsto x+\eps_Nk$ can be placed on a zero of $\sh$ while $x$ is
not, that supremum is $+\infty$ for every $k\neq0$ \emph{in both sectors}. The step from (209)
to (210) is then carried entirely by the sentence ``\emph{as long as $N$ is sufficiently large,
we expect essentially the same behavior for $k\neq0$ because of the uniform continuity of
$\sh$}'', which is the invalid step of Finding~\ref{f:lemma45} in supremum form.

\medskip\noindent
Measured (\texttt{os\_check2.py} [7], \texttt{os\_check3.py} [10]), $\mathrm{db}4$, $\delta=2$:
\begin{center}\small
\begin{tabular}{lcccc}
\toprule
 & $N=6$ & $N=8$ & $N=10$ & $N=12$\\
\midrule
$\mathrm{Error}^2$, Ramond, $k=1$ & $+\infty$ & $+\infty$ & $\ge3.9\cdot10^{12}$ & --- \\
$\mathrm{Error}^2$, NS, $k=1$     & $1.21\cdot10^{-7}$ & $5.61\cdot10^{-10}$ & $2.35\cdot10^{-12}$ & $9.44\cdot10^{-15}$\\
successive ratios (NS)            & \multicolumn{4}{l}{$215,\ 239,\ 249$ \quad vs.\ $2^{2\delta\cdot 2}=256$ predicted by (210)}\\
\bottomrule
\end{tabular}
\end{center}
So the \emph{conclusion} (210) is numerically correct in the Neveu--Schwarz sector, false in the
Ramond sector, and derived by an invalid argument in both. Since the paper's own footnote~15
already flags that the Ramond sector needs extra care with $\chi_{\Gamma_N}$, this is a natural
place for the fix.

\paragraph{What is unaffected --- and it is the part that matters downstream.}
Remark 4.21 / eq.~(212), the $k=0$ time-evolution bound
\[
  \big\|\big((\tilde\ell^{(N)}_{\pm,0}-\ell_{\pm,0})_{\pm\pm}R^M_\infty(\eta_p)\big\|^2_\infty
  \le \eps_N^{2\delta}\Big(\sup_{x}\tfrac{|\mathrm{sinc}(x)-1|^2}{|x|^{2\delta}}\Big)
  \|\eta\|_M^2(\tfrac L\pi)^2\!\!\sum_{n\in\Gamma_{M,\pm}}\!\!\theta(\pm n)|n|^{2(1+\delta)},
\]
comes from the \emph{momentum-cutoff} route (201) and contains no $\sh$ at all: the sum is
finite (over the $M$-th Brillouin zone) and no regularity of $s$ enters. I confirmed
$\sup_x|\mathrm{sinc}(x)-1|^2/|x|^4=1/36$, attained as $x\to0$ (it follows from
$1-x^2/6\le\mathrm{sinc}(x)\le1$). \emph{The advertised $\sim2^{-2N}$ error scaling for the
time evolution is sound and unconditional.} This is the estimate that feeds the companion
quantum-simulation paper; the $k\neq0$ budget from (210) is not.

\section{Finding 3: what ``sufficiently regular'' actually costs}\label{sec:reg}

\begin{finding}
The hypothesis ``$s\in C^\alpha(\mathbb R)$ sufficiently regular'' is never quantified, and
carries \emph{three different} thresholds. In terms of the Sobolev exponent $\sigma_K$ of the
Daubechies-$K$ scaling function:
\begin{enumerate}[label=(\alph*),leftmargin=1.9em,itemsep=1pt]
\item $\DW\subset\mathcal D(\ell_{\pm,k})$ (needed for $\DW$ to be a core at all, p.~44) and
$\ell^2$-domination in Lemma 4.5 both require $\sigma_K>1$;
\item the two Cauchy--Schwarz splittings in the proof of Lemma 4.15 require $\delta>3/2$
(second factor, since $\|\tilde\ell^{(N)}_{\pm,k}R^K_\infty(e_l)\|$ grows linearly in $k$) and
then $\sigma_K>\delta$ (first factor, since $\hat X_{\mathrm{per}}$ is periodic and does not
decay) --- hence $\sigma_K>3/2$;
\item finiteness of $\|\sh(\eps_M(\,\cdot+\tfrac k2))\|_{h^{1+\delta}}$ in (207) requires
$\sigma_K>1+\delta$; for the $\delta=2$ used to advertise the $2^{-2\delta N}$ rate this is
$\sigma_K>3$.
\end{enumerate}
\end{finding}

\noindent
I constructed the Daubechies filters from scratch (spectral factorisation of
$\cos^{2K}(\xi/2)P_K(\sin^2(\xi/2))$; orthonormality verified to $<6\cdot10^{-15}$) and measured
$\sigma_K$ from dyadic-block $L^2$ decay of $\sh$; the values reproduce the standard
Villemoes/Daubechies table to 3--4 digits (\texttt{os\_check4.py}).

\begin{center}\small
\begin{tabular}{lccccccccccc}
\toprule
$K$ & 2 & 3 & 4 & 5 & 6 & 7 & 8 & 9 & 10 & 11 & 12\\
\midrule
$\sigma_K$ (measured) & 1.002 & 1.416 & 1.776 & 2.098 & 2.390 & 2.661 & 2.917 & \textbf{3.165} & 3.406 & 3.641 & 3.873\\
$K-\bar K$ (sup-envelope) & 1.336 & 1.634 & 1.910 & 2.174 & 2.431 & 2.681 & 2.927 & 3.170 & 3.409 & 3.646 & 3.881\\
\midrule
(a) $\sigma_K>1$      & \color{red}$\approx$ & \checkmark & \checkmark & \checkmark & \checkmark & \checkmark & \checkmark & \checkmark & \checkmark & \checkmark & \checkmark\\
(b) $\sigma_K>3/2$    & \color{red}$\times$ & \color{red}$\times$ & \checkmark & \checkmark & \checkmark & \checkmark & \checkmark & \checkmark & \checkmark & \checkmark & \checkmark\\
(c) $\sigma_K>3$      & \color{red}$\times$ & \color{red}$\times$ & \color{red}$\times$ & \color{red}$\times$ & \color{red}$\times$ & \color{red}$\times$ & \color{red}$\times$ & \checkmark & \checkmark & \checkmark & \checkmark\\
\bottomrule
\end{tabular}
\end{center}

\noindent
Three consequences worth stating in a revision:
\begin{itemize}[leftmargin=1.4em,itemsep=1pt]
\item $\mathrm{db}2$ is \emph{exactly borderline} ($\sigma_2=1$, so $s\notin H^1$): the wavelet
core $\DW$ is then not contained in $\mathcal D(\ell_{\pm,k})$ and Lemma 4.5 has no content.
The honest minimum for Section 4.2.4 is $K\ge3$, and $K\ge4$ once Lemma 4.15 is included.
\item The advertised $\delta=2$ rate in (210) requires $K\ge9$ --- \emph{not} $K\ge2$, and not
the ``$K\ge6$ suffices'' recorded in the earlier verification note in
\texttt{osborne-stottmeister2023/}. Since $\mathrm{supp}({}_K s)=[0,2K-1]$, Remark 4.20's own
constraint $\eps_M(2K-1)\sim L$ then forces a correspondingly coarser $M$: the regularity
needed for the rate and the localisation needed for the qubit count pull against each other,
and the paper does not resolve the trade-off.
\item The relevant exponent differs between the two uses: the pointwise bound (109) supplies
$K-\bar K$, whereas the $\ell^2$-sums actually need the (larger) Sobolev exponent. Using (109)
to check the sums, as the paper does, is lossy by roughly $0.4$--$0.5$ in $K$.
\end{itemize}

\section{Finding 4: Remark 4.14 does not follow from Lemma 4.5}

\noindent
Remark 4.14 states that in the NS sector the Hilbert--Schmidt norms in (201) vanish for
$k=0,\pm\pi/L$ (the Moebius generators), and concludes that ``\emph{a simple extension of
Lemma 4.5 is sufficient to prove Theorem 4.11 in this case using the wavelet renormalization
group}.''

The vanishing is correct, though for a subtler reason than emptiness of the index set: for
$k=-\pi/L$ and the $+$ chirality the constraint $\theta(n-k)\theta(-n)$ admits
$n=-\pi/2L\in\Gamma_{\infty,-}$, and the summand vanishes there because $n-\tfrac k2=0$ kills
both $\sin(\eps_N(n-\tfrac k2))$ and $(n-\tfrac k2)$.

\begin{finding}
The conclusion does not follow. Estimate (200) also requires the \emph{diagonal} blocks
$(\tilde o^{(N)}-o)_{\pm\pm}=S_\pm(\tilde o^{(N)}-o)S_\pm$ applied to $R^M_\infty(\eta)$, i.e.\
convergence on the Hardy-projected vectors $P^\mp\xi$. But $P^\pm$ does not map $\DW$ into
$\DW$.
\end{finding}

\noindent
$\DW=\bigcup_NR^N_\infty(\hs_{N,\pm})=\bigcup_NV_N$ is the union of the multiresolution spaces;
a $V_M$ element has $\hat\xi(k)=\eps_M^{1/2}\sh(\eps_Mk)\hat\eta_{\mathrm{per}}(k)$ with
$\hat\eta_{\mathrm{per}}$ periodic, so it cannot be supported on positive frequencies alone
without vanishing. Measured relative distance of $P^+s^{(\eps_M)}$ from $V_J$
(\texttt{os\_check3.py} [11], $\mathrm{db}4$, $M=3$): $0.353,\,0.080,\,0.0074,\,6.4\cdot10^{-4},
\,5.5\cdot10^{-5}$ for $J=3,4,6,8,10$ --- convergent to $0$, as it must be by density, but never
$0$.

\begin{claim}[repair]
$\ell_{\pm,k}$ and $\tilde\ell^{(N)}_{\pm,k}$ are bounded $h^1(\Gi)\to h^0(\Gi)$ \emph{uniformly
in $N$} (their symbols are $O(|m|)$, and $(R^N_\infty)^*$ is bounded $h^1\to h^1(\Gamma_N)$ by
the decay of $\sh$). Since $\DW$ is $h^1$-dense (the paper's own graph-norm argument, p.~44)
and $P^\pm$ is a Fourier multiplier preserving $h^1$, strong convergence extends from $\DW$ to
all of $h^1\ni P^\mp\xi$. This closes the remark, but it is not ``a simple extension of Lemma 4.5''.
\end{claim}

\noindent
This also isolates an unremarked structural reason why the momentum-cutoff RG is needed: in
that setting $P^\pm$ is diagonal in the plane-wave basis, so $\Dstd$ \emph{is} invariant and
(200) applies verbatim. The paper motivates the switch by the Hilbert--Schmidt condition alone.

\section{Finding 5: a normalisation inconsistency in Definition 4.2}

\noindent
Two incompatible normalisations of the one-particle Virasoro generator are in circulation:
\[
\text{(A) eq.\ (143)/(165):}\quad \ell_k(l)=-(l+\tfrac k2)p_-,
\qquad\qquad
\text{(B) Def.\ 4.2 (181) + (199):}\quad \ell_{\pm,k}e_n=\pm\tfrac L\pi(n\mp\tfrac k2)e_{n\mp k}.
\]
(B) is internally consistent --- (181), (184), (198), (199) and the estimates (200)/(201) all use
it, and $\tilde\ell^{(N)}_{\pm,k}\to\ell_{\pm,k}$ holds in (B) --- but it is (A) that is
compatible with the bridge formula $L_{\pm,k}=\frac L\pi\,dF_S(\ell_{\pm,k})$ used in (142),
(145), (167), (170) and (194). The two differ by exactly $L/\pi$.

\begin{finding}
Read literally with Definition 4.2, the commutator (170) acquires the structure constant
$(L/\pi)^2(k-k')$ instead of $\frac L\pi(k-k')$, and its central term becomes
$(L/\pi)^2\cdot\frac1{12}n(n^2-1)$ with $n=\frac L\pi k$, i.e.\ $c=(L/\pi)^2$ rather than $c=1$.
The same $L/\pi$ propagates into (173), (177), (178) and (194).
\end{finding}

\noindent
Direct evaluation of the Schwinger cocycle $c_{P^\mp}(\ell_{+,k},\ell_{+,-k})$ over the finite
strip $0<|l|<|k|$ (\texttt{os\_check5.py}), target $\frac1{12}n(n^2-1)$:
\begin{center}\small
\begin{tabular}{cccc}
\toprule
$L$ & $n$ & $(L/\pi)^2c_S$ with (A) & $(L/\pi)^2c_S$ with (B)\\
\midrule
$\pi$   & 3 & 2.000 & 2.000\\
$2\pi$  & 3 & 2.000 & \color{red}8.000\\
$5$     & 3 & 2.000 & \color{red}5.066\\
$2\pi$  & 5 & 10.000 & \color{red}40.000\\
\bottomrule
\end{tabular}
\end{center}
The two agree iff $L=\pi$, which the paper nowhere assumes ($L$ is free throughout;
$\eps_NL_N=L$). \emph{No convergence statement is affected} --- the factor is common to
approximant and limit --- but the identification of the limit with the $c=1$ (resp.\ $c=\frac12$)
Virasoro generators requires normalisation (A), so Definition 4.2 should read
$\ell_{\pm,k}=\pm\frac{1}{2L}\sum_l(l\mp\frac k2)\,e_{l\mp k}\otimes e_l$ (equivalently, drop
$L/\pi$ from (199)).

\section{Finding 6: Theorem 6.3 rests on two unproved assertions}

\noindent
The telescoping estimates (236)/(237) reduce the statement to applying Theorem 4.16 to the
vectors $\prod_{p<q}L_\pm(X_p)\Omega_0$. The proof then closes with:
\begin{quote}\small
``\emph{First, $\|\prod_{p=1}^n\!:\!(\pi_\pm(\alpha_\infty^N(L^{(N)}_\pm(X_{N,p}))))\!:\Omega_0\|$
can be bounded uniformly in $N$ \dots\ Second, $\prod_{p=1}^nL_\pm(X_p)\Omega_0$ is given by a
convergent sum \dots\ of Fock-space vectors with finite particle number in momentum space.}''
\end{quote}

\begin{finding}
These two ``observations'' are exactly the decisive ingredients and neither is proved.
\end{finding}

\noindent
They are not decorative. Theorem 4.16 gives strong convergence only on
$\mathfrak F^{\mathrm{alg}}_{\mathrm a}(\Dstd)$, and the intermediate vectors are \emph{not} in
that domain. Note the contrast with the unsmeared case: for fixed $k$ the off-diagonal block
$(\ell_{\pm,k})_{\pm\mp}$ is supported on the finite strip $0<|n|<|k|$, hence \emph{finite rank},
so $L_{\pm,k}\Omega_0$ is a finite combination of plane-wave pairs. After smearing,
$L_\pm(X)=\frac1{2L}\sum_k\hat X_k L_{\pm,k}$ has an off-diagonal block that is
Hilbert--Schmidt but of infinite rank, so $L_\pm(X)\Omega_0=a^\dagger G_{+-}a^\dagger\Omega_0$
is an infinite two-particle superposition, outside
$\mathfrak F^{\mathrm{alg}}_{\mathrm a}(\Dstd)$ and of unbounded $L_{\pm,0}$-energy.

What is needed is an \emph{$N$-uniform energy bound}
$\|:\!\pi_\pm(\alpha^N_\infty(L^{(N)}_\pm(X)))\!:\Psi\|\le C_X\|(1+L_{\pm,0})\Psi\|$, the lattice
analogue of the Buchholz--Schulz-Mirbach / Carpi--Weiner energy bounds for smeared Virasoro
fields. That is a theorem in its own right, and it is the main rigor debt of Section 6. The
same gap does not affect Theorem 6.1, where the insertions are bounded elements of the CAR
algebra and the argument is complete.

\medskip\noindent
Two smaller items in the same vicinity. Theorem 4.11 is stated with no proof at all (it is
offered as an immediate consequence of Lemma 4.10 and (200)); and the assertion that
$\mathfrak F^{\mathrm{alg}}_{\mathrm a}(\Dstd)$ contains a total set of analytic vectors for
$L_{\pm,k}$ is justified only by a reference to Corollary 4.3 and (200).

That assertion does not follow, and it cannot be had by iterating (25). In $\pi_\pm$ the pair
terms $a^\dagger G_{+-}a^\dagger$ and $aG_{-+}a$ \emph{raise} the particle number, so at step
$j$ the estimate is applied at level $n+2j$ rather than at level $n$, and iterating (25)
together with (187) produces a factor $(j!)^2$ after $j$ steps; the exponential series then
diverges for every $t>0$.\footnote{The first version of this note asserted the opposite,
deriving a radius $t<1/(2\|G_{+-}\|_2)$ from
$\|(a^\dagger G_{+-}a^\dagger)^jP_{\le n}\|\le 2^jj!\,j^{n/2}\|G_{+-}\|_2^{\,j}$. That bound
holds the particle number fixed at $n$ and is wrong.} The statement itself is presumably true,
being the Goodman--Wallach/Carpi--Weiner analyticity of finite-energy vectors, but it does not
follow from the estimates available here.

What repairs Theorem 4.11 is therefore not an analytic-vector argument at all, but Nelson's
\emph{commutator} theorem [Reed--Simon X.37] applied with $N=1+L_{\pm,0}$: its first hypothesis
is exactly the uniform energy bound demanded above, and its second follows from the Virasoro
relation $[L_{\pm,0},L_{\pm,k}]=-\tfrac L\pi kL_{\pm,k}$ together with the form bound. The
conclusion is stronger than the one claimed --- essential self-adjointness on \emph{every} core
for $1+L_{\pm,0}$, in particular on $\mathcal D(L_{\pm,0})$.

\section{Finding 7: the locality claim for Theorem 4.16}

\noindent
The preamble to Lemma 4.15 states that combining the two renormalisation groups exploits
convergence in $\pi_\pm$ ``\emph{while, at the same time, preserving localization in real space
in the sense of Proposition 3.6}.''

\begin{finding}
This is not established. Proposition 3.6 is a statement about the \emph{wavelet} RG. In
Theorem 4.16 the fermion algebra is transported by the \emph{momentum-cutoff} maps
$\alpha^N_\infty$, which by the paper's own eq.~(112) send lattice-localised operators to
operators that are not localised. The wavelet RG enters only through the smearing functions
$S^M_N(X)$, and inspection of the proof of Lemma 4.15 shows its role there is summability of
the $k$-sum (the weights $(1+|\eps_Mk|)^{\pm2\delta}$), not locality.
\end{finding}

\noindent
Compounding this, the approximants used in Theorem 4.16 are the \emph{modified} ones (197),
whose real-space form (159)--(162) is a convolution with $\widehat{\chi_{\Gamma_N}}$ --- a
Dirichlet kernel with only $O(1/|x|)$ decay. What is true, and worth saying instead, is that
the smearing function is wavelet-localised at scale $M$; the observable is quasi-local only up
to Dirichlet tails. The distinction matters for the quantum-simulation reading of these results,
where locality of the simulated operator is the whole point.

\section{Finding 8: minor corrigenda}

\begin{enumerate}[leftmargin=1.7em,itemsep=2pt]
\item \textbf{(25), p.~15.} ``$\|aG_{+-}aP_{\le n}\|\le n\|G_{+-}\|_2$'' should read
$\|aG_{-+}aP_{\le n}\|\le n\|G_{-+}\|_2$ --- which is how it is used in (200) and in (22)/(24).
\item \textbf{Lemma 3.7, p.~27.} The display asserts ``$=$'' where ``$\le$'' is meant. Also
$\max\{e^C,\pi^{-K_j}\}=e^C$ identically, since $q_j\ge1$ forces $K_j\ge0$; and case 3 of the
proof claims $2^{MK_j}\le\pi^{-K_j}(1+|l|)^{K_j}$, which fails throughout
$2^M<|l|<\pi2^M-1$. The correct step is simply $2^{MK_j}\le(1+|l|)^{K_j}$, i.e.\ drop
$\pi^{-K_j}$.
\item \textbf{(204)/(205), p.~53.} The factorisation divides by $(m+\tfrac k2)$, which vanishes
at $m=-k/2$ whenever $k\in\frac{2\pi}{L}\mathbb Z$ (removable, since $\sin$ vanishes there too).
The Sobolev norm in (205)--(207) is written with $\sh(\eps_M(\,\cdot+\tfrac k2))$ while (204)
carries $\sh(\eps_M(\,\cdot+k))$.
\item \textbf{(213), p.~56.} From $\binom{2j}{j}^{1/2}\le2^j$ and
$\prod_{i=0}^{j-1}(a+ib)\le a^j(2j-1)!!\le a^j2^jj!$ with $a=|m|+\tfrac12|k|$, $b=|k|\le2a$, the
bound is $(4a)^jj!$, not $(2a)^jj!$; (213) should then carry $\frac{2L}{\pi}(|m|+\frac12|k|)t$
in place of $\frac L\pi(|m|+\frac12|k|)t$. The conclusion (a positive radius of convergence) is
unchanged.
\item \textbf{§4.2.2.} The lattice Virasoro relation (146) and the $O(\eps_N^2)$ remainders
(147) are computed for the \emph{unmodified} approximants; the modified ones (160)/(162)
carrying $\chi_{\Gamma_N}$ --- which are the ones used in Theorem 4.11 and 4.16 --- are only
argued to agree asymptotically. Their finite-$N$ algebra is not checked. Harmless for the
theorems, but the sentence ``we recover the original Koo--Saleur formulas for $N\to\infty$''
is doing more work than it looks.
\end{enumerate}

\section*{Verdict}

\noindent
Theorems A, B and C (= 4.11, 6.1, 6.3, with 4.6, 4.16, 5.1) are, as far as I can tell, true, and
the architecture --- one-particle $\to$ Fock $\to$ correlation functions, with two
renormalisation groups doing two different jobs --- is sound and, in my view, the right way to
do this. But the earlier ``no logical gaps'' verdict does not survive contact with the details.
Ranked by what a referee should have caught:

\begin{enumerate}[leftmargin=1.7em,itemsep=2pt]
\item \textbf{The $\sh$-quotient (Findings 1--2).} An invalid step, used three times, that is
\emph{false} rather than merely unjustified, and that makes the advertised $k\neq0$ error
estimate infinite in the Ramond sector. Lemma 4.5 is repairable in a paragraph using the
paper's own Lemma 3.7; (210) needs genuine rewriting for $k\neq0$.
\item \textbf{Theorem 6.3 (Finding 6).} Two unproved assertions carry the proof, and one of them
(an $N$-uniform energy bound) is a real theorem.
\item \textbf{Regularity (Finding 3).} ``Sufficiently regular'' hides $K\ge3$, $K\ge4$ and
$K\ge9$ in three different places. The $K\ge9$ is the interesting one: the paper's headline
practical claim --- explicit error bounds for quantum simulation --- is stated at $\delta=2$,
which no Daubechies wavelet below order $9$ supports along the wavelet route.
\item \textbf{Findings 4, 5, 7.} Each is a paragraph-sized fix, but 5 would read as a factual
error about the central charge to anyone who takes Definition 4.2 at face value.
\end{enumerate}

\noindent
Net: the paper is right, and it is under-proved in a way that a careful revision could fix in
about three pages. The one thing I would not simply patch is (210): the honest statement is
that the wavelet route gives an explicit rate only for $k=0$, and that the clean, unconditional
error bound available today is the momentum-cutoff estimate (212).

\appendix
\section{Numerical protocol}

\noindent
All checks are in \texttt{lattice\_cft/numerics/} (pure \texttt{numpy}, no external wavelet
library). Conventions: $L=\pi$, so $\Gamma_{\infty,+}=\mathbb Z$,
$\Gamma_{\infty,-}=\mathbb Z+\tfrac12$, $\eps_N=2^{-N}\pi$;
$m_0(\xi)=2^{-1/2}\sum_nh_ne^{-in\xi}$, $\sh(\xi)=\prod_{j\ge1}m_0(2^{-j}\xi)$ truncated
adaptively ($J=\lceil\log_2\max|\xi|\rceil+32$, the tail being $1+O(2^{-J}|\xi|)$).

\begin{itemize}[leftmargin=1.4em,itemsep=1pt]
\item \texttt{os\_check.py} --- filter sanity ($\sum h_n=\sqrt2$, $\langle h,h(\cdot-2k)\rangle=\delta$
to $<10^{-12}$), order of the zeros of $\sh$ at $2\pi n$, the $\sh$-quotient in both sectors.
\item \texttt{os\_check2.py} --- cross-check of the product form against the quotient form of
$f^{(N)}_k$ ($<2\cdot10^{-14}$ where both are defined); failure of the stated dominating
function; $\ell^2$-convergence of $f^{(N)}_k\to f_k$; $2^{-KN}$ decay of the resonance spikes;
$\mathrm{Error}^2$ in the Ramond sector.
\item \texttt{os\_check3.py} --- decay exponents; existence of a valid dominating function;
$\mathrm{Error}^2$ in the NS sector; distance of $P^+\!s^{(\eps_M)}$ from $V_J$.
\item \texttt{os\_check4.py} --- Daubechies filters for $K=2,\dots,12$ by spectral factorisation
(orthonormality $<6\cdot10^{-15}$); Sobolev exponents $\sigma_K$ from dyadic-block $L^2$ decay,
reproducing the literature values ($\sigma_2=1.000$, $\sigma_3=1.415$, $\sigma_4=1.776$,
$\sigma_6=2.386$, $\sigma_8=2.917$) to 3--4 digits.
\item \texttt{os\_check5.py} --- Schwinger cocycle over the finite strip, both normalisations.
\end{itemize}

\end{document}
